\origpage{73}{GRAPH MINORS. XI}
\noindent it by the use of ``tangles,'' developed in earlier papers. Roughly, a tangle of order $\theta$ in a graph can be thought of as a highly coherent area of the graph, which clearly resides on one side or the other of every separation of order $< \theta$. If we have a graph drawn on the sphere with a tangle of order $\ge \theta$ then for any simple closed curve $C$ which meets the drawing $< \theta$ times (and which meets the drawing only in vertices), the two discs it bounds define a separation of the graph of order $< \theta$, and we may define the ``inside'' of $C$ to be whichever of these discs does not ``include'' the tangle. We find that, with this definition of ``inside,'' we can handle drawings on the sphere just as easily as drawings on higher surfaces. Conveniently, there is a complementary result that, if a graph has a $\theta$-representative drawing on a connected surface not the sphere, this defines a tangle in the graph of order $\theta$; and this tangle resides ``outside'' every null-homotopic curve which meets the drawing $< \theta$ times, just as in the sphere case. This permits us to completely unify the sphere and non-sphere cases.

Our first objective in this paper is to prove that a graph with a $\theta$-representative drawing on a connected surface not the sphere contains an associated tangle of order $\theta$. This occupies Sections 5 and 6. (The earlier sections are concerned with definitions.) Secondly, given a graph on a surface (highly representative if the surface is not a sphere, and with a designated high order tangle if the surface is a sphere, so that in either case we have a definition of ``inside``), there is associated a natural distance function on the surface, and in Sections 7-9 we study some of its properties. In particular, it turns out that for any point $c$ of the surface, the set of all points within distance $\theta'$ of $c$ (for any $\theta' < \theta$) is simply-connected, and thus every point has a ``large'' planar neighborhood. This will be useful in subsequent papers.

We are concerned with simple closed curves in the surface which meet the drawing $\Gamma$ only in vertices, and only a small number of times. There is an associated graph (the ``radial drawing'' of $\Gamma$) in which such curves correspond to circuits, and it is most convenient to work in terms of this rather than directly in terms of $\Gamma$.

One can also draw hypergraphs on surfaces, and define tangles in hypergraphs, and it turns out that all our results generalize to hypergraphs. (The radial drawing of a drawing of a hypergraph is still a graph, not a hypergraph, and so this generalization to hypergraphs makes very little difference.) Thus we formulate the present results in terms of hypergraphs.

\sectionhead{2. Hypergraphs and Tangles}

A hypergraph $G$ consists of a set $V(G)$ of vertices, a set $E(G)$ of edges, and an incidence relation between them. All hypergraphs in this paper are
